\section{Appendix: Air Space}
\label{sec:appendix_airspace}

This appendix describes the JSON format consumed by the 'Import Air Space...' action of the Data Context Sectors editor (see \nameref{sec:ui_configure_sectors}). The format differs from the COMPASS native sector JSON (covered by 'Import JSON...' / 'Export JSON...') in that it carries altitude bounds per sector and allows a single level of nesting (a top-level entry can act as a layer that groups sub-sectors). \\

\subsection{Top Level}

The file consists of a single object with one key, which contains an array of sectors:

\begin{lstlisting}
{
    "air_space_sector_definition": {
        "air_space_sectors": [
            ...
        ]
    }
}
\end{lstlisting}

Each element of the 'air\_space\_sectors' array is a sector object as described below. All numeric values are encoded as JSON strings. \\

\subsection{Sector Object}

Each sector object provides the following fields:

\begin{itemize}
\item name: Sector name (string)
\item own\_volume: '1' if the sector defines its own polygon volume, '0' otherwise (string)
\item used\_for\_checking: '1' if the sector is intended to be used as a checking volume, '0' otherwise (string)
\item own\_height\_min: Lower altitude bound in feet (string)
\item own\_height\_max: Upper altitude bound in feet (string)
\item geographic\_points: Array of polygon vertices, each carrying an integer 'index' (string), 'latitude' (WGS-84 degrees, string) and 'longitude' (WGS-84 degrees, string). Vertices are ordered by ascending index value.
\item air\_space\_sectors: Optional array of nested sector objects, turning the parent into a layer that groups them. Only a single level of nesting is supported.
\end{itemize}
\ \\

Please \textbf{note} that all numeric fields are stored as JSON strings, not as JSON numbers. Altitude bounds are given in feet. \\

\subsection{Flat Example}

The following example defines a single sector 'Upper' with four polygon vertices and an altitude band between 28500 ft and 60000 ft:

\begin{lstlisting}[basicstyle=\small\ttfamily]
{
    "air_space_sector_definition": {
        "air_space_sectors": [
            {
                "name": "Upper",
                "own_volume": "1",
                "used_for_checking": "0",
                "own_height_min": "28500.000000",
                "own_height_max": "60000.000000",
                "geographic_points": [
                    { "index": "0", "latitude": "47.025600",
                      "longitude":  "8.636990" },
                    { "index": "1", "latitude": "47.191700",
                      "longitude":  "8.560150" },
                    { "index": "2", "latitude": "48.199100",
                      "longitude":  "9.366450" },
                    { "index": "3", "latitude": "45.708300",
                      "longitude": "14.571000" }
                ]
            }
        ]
    }
}
\end{lstlisting}

\subsection{Nested Example}

A top-level entry may carry an 'air\_space\_sectors' array of its own, in which case it is treated as a layer and its children become the sectors of that layer. In the example below, the layer 'TMAs' groups two terminal sectors: \\

\begin{lstlisting}[basicstyle=\small\ttfamily]
{
    "air_space_sector_definition": {
        "air_space_sectors": [
            {
                "name": "TMAs",
                "own_volume": "1",
                "used_for_checking": "0",
                "own_height_min": "0.000000",
                "own_height_max": "24500.000000",
                "geographic_points": [ ... ],
                "air_space_sectors": [
                    {
                        "name": "TMAVienna",
                        "own_volume": "1",
                        "used_for_checking": "0",
                        "own_height_min": "0.000000",
                        "own_height_max": "24500.000000",
                        "geographic_points": [
                            { "index": "0", "latitude": "48.110000",
                              "longitude": "16.570000" },
                            { "index": "1", "latitude": "48.250000",
                              "longitude": "16.720000" },
                            { "index": "2", "latitude": "48.300000",
                              "longitude": "16.350000" }
                        ]
                    },
                    {
                        "name": "TMAGraz",
                        "own_volume": "1",
                        "used_for_checking": "0",
                        "own_height_min": "0.000000",
                        "own_height_max": "16500.000000",
                        "geographic_points": [
                            { "index": "0", "latitude": "47.000000",
                              "longitude": "15.430000" },
                            { "index": "1", "latitude": "47.080000",
                              "longitude": "15.520000" },
                            { "index": "2", "latitude": "47.020000",
                              "longitude": "15.290000" }
                        ]
                    }
                ]
            }
        ]
    }
}
\end{lstlisting}

When such a file is loaded via 'Import Air Space...', each top-level sector and each nested sector appears as a row in the import dialog's sector tree, where individual sectors can be (de)selected before being added to the active Data Context under the user-chosen layer name. \\
